% ทฤษฎีและงานวิจัยที่เกี่ยวข้อง ภาษาไทย ใช้ใน main_short_th.tex (ฉบับอังกฤษอยู่ใน 02_related_work.tex);
% แก้ให้ตรงกันทั้งสองฉบับ. ไม่ใช้ \emph เพราะฟอนต์ไทยไม่มีตัวเอียง.
\section{ทฤษฎีและงานวิจัยที่เกี่ยวข้อง}\label{sec:related}

บทนี้อธิบายทฤษฎีที่งานวิจัยนี้อาศัย ได้แก่ เหตุใดภาพเดียวจึงบอกความลึกเป็นเมตรไม่ได้
โมเดลที่เรานำมาปรับจูนมีโครงสร้างและฝึกมาอย่างไร และความลึกรายภาพกลายเป็นแบบจำลองห้องได้อย่างไร
พร้อมทั้งวางตำแหน่งของงานนี้เทียบกับงานก่อนหน้าในแต่ละประเด็น

\subsection{การประมาณความลึกจากภาพเดี่ยว}
ตามแบบจำลองกล้องรูเข็ม (pinhole camera)~\cite{hartley} ขนาดของวัตถุในภาพแปรผันตามขนาดจริงของวัตถุ
หารด้วยระยะห่าง แล้วคูณด้วยความยาวโฟกัสของเลนส์ ห้องที่ใหญ่เป็นสองเท่าเมื่อมองจากระยะไกลเป็นสองเท่า
จึงให้ภาพเหมือนกันทุกประการ เรขาคณิตเพียงอย่างเดียวจึงหาความลึกเป็นเมตรจากภาพเดียวไม่ได้
โครงข่ายประสาทเทียมทำได้เพียงอนุมานจากสิ่งที่เรียนรู้มา เช่น ความสูงปกติของประตูหรือโต๊ะ
และแม้แต่ข้อมูลเหล่านี้ก็ขึ้นกับความยาวโฟกัส เพราะวัตถุเดียวกันที่ระยะเดียวกันจะดูใหญ่ขึ้นเมื่อใช้เลนส์ที่ยาวขึ้น
Eigen และคณะ~\cite{eigen14} เป็นงานแรก ๆ ที่ใช้โครงข่ายเชิงลึกกับปัญหานี้
และเพราะความกำกวมดังกล่าว จึงวัดความคลาดเคลื่อนบนลอการิทึมของความลึก
ในรูปแบบที่แยกความคลาดเคลื่อนของสเกลโดยรวมของฉากออกจากความคลาดเคลื่อนของรูปทรง
ฟังก์ชันสูญเสีย (loss) แบบไม่ขึ้นกับสเกลบนลอการิทึมนี้ในรูปแบบถ่วงน้ำหนัก เรียกว่า SiLog
ใช้ฝึกทั้ง ZoeDepth~\cite{zoedepth} และโมเดลที่ปรับจูนในงานนี้

โมเดลแบ่งได้เป็นสองกลุ่มตามวิธีจัดการกับสเกลที่หายไป กลุ่มแรกคือโมเดลความลึกสัมพัทธ์ (relative) ซึ่งไม่พยายามหาสเกล
MiDaS~\cite{midas} ฝึกโครงข่ายเดียวด้วยชุดข้อมูลหลายชุดที่วัดความลึกด้วยเซนเซอร์และหน่วยต่างกัน
ซึ่งทำได้ก็ต่อเมื่อฟังก์ชันสูญเสียไม่สนใจสเกลและค่าเลื่อน (offset) ของแต่ละภาพ
ผลลัพธ์จึงเป็นส่วนกลับของความลึก (disparity) ที่ถูกต้องเฉพาะเมื่อรู้สเกลและค่าเลื่อนของภาพนั้น
ต้องหาตัวเลขสองตัวนี้จากข้อมูลอื่นก่อนจึงจะได้ความลึกเป็นเมตร Depth Anything~\cite{da1}
และ Depth Anything~V2~\cite{da2} (\mbox{DA-v2}) ให้ผลลัพธ์แบบเดียวกัน แต่ฝึกด้วยภาพจำนวนมากกว่ามาก
กลุ่มที่สองคือโมเดลความลึกแบบเมตริก (metric) ซึ่งให้ความลึกเป็นเมตรโดยตรงโดยเรียนรู้สเกลจากข้อมูลที่มีค่าความลึกกำกับ
ZoeDepth~\cite{zoedepth} ต่อส่วนทำนายแบบเมตริกที่ปรับจูนด้วย NYU Depth~v2~\cite{nyuv2} และ KITTI~\cite{kitti}
เข้ากับโมเดลความลึกสัมพัทธ์ ส่วน \mbox{DA-v2} มีรุ่นเมตริกที่ปรับจูนจากรุ่นสัมพัทธ์ของตนเอง
โดยรุ่นในอาคาร (\mbox{DA-v2} Metric-Indoor) ปรับจูนด้วยภาพห้องที่เรนเดอร์จากชุดข้อมูล Hypersim~\cite{hypersim}
อีกแนวทางหนึ่งจัดการกับความยาวโฟกัสโดยตรง Metric3D~\cite{metric3d} แปลงทุกภาพให้เป็นภาพจากกล้องมาตรฐานที่มีความยาวโฟกัสคงที่
UniDepth~\cite{unidepth} ทำนายพารามิเตอร์ของกล้องไปพร้อมกับความลึก และ Depth Pro~\cite{depthpro}
ประมาณความยาวโฟกัสเองเมื่อไม่ได้รับค่านี้มา งานเหล่านี้รายงานผลเป็นตัวชี้วัดรายภาพบนชุดทดสอบมาตรฐาน
เช่น NYU Depth~v2 และ KITTI ตัวชี้วัดดังกล่าวเฉลี่ยจากภาพที่เป็นอิสระต่อกัน จึงไม่บอกว่าสเกลของโมเดลคงที่หรือไม่
ตลอดเฟรมของการสแกนครั้งเดียว ซึ่งเป็นสิ่งสำคัญเมื่อนำเฟรมทั้งหมดมารวมเป็นห้องเดียว

\subsection{โครงสร้างและการฝึกของ Depth Anything V2}
\mbox{DA-v2} ประกอบด้วยส่วนเข้ารหัสภาพ (encoder) และส่วนถอดรหัสความลึก (decoder)
ส่วนเข้ารหัสคือ DINOv2~\cite{dinov2} ซึ่งเป็น vision transformer (ViT)~\cite{vit}
โดยแบ่งภาพเป็นชิ้นสี่เหลี่ยมเล็ก ๆ (patch) แล้วให้แต่ละชิ้นดึงข้อมูลจากชิ้นอื่นทุกชิ้นผ่านกลไกที่เรียกว่า self-attention
DINOv2 ผ่านการฝึกล่วงหน้าโดยไม่ใช้ข้อมูลกำกับ บนชุดภาพขนาดใหญ่ที่คัดสรรแล้ว
เพื่อให้คุณลักษณะ (feature) ที่ได้ใช้กับงานได้หลากหลาย ส่วนถอดรหัสคือ Dense Prediction Transformer (DPT)~\cite{dpt}
ซึ่งนำคุณลักษณะจากหลายชั้นของส่วนเข้ารหัสมาประกอบกลับเป็นแผนที่คล้ายภาพที่หลายความละเอียด
แล้วรวมเป็นแผนที่ความลึกแบบหนาแน่นเพียงแผนที่เดียว โค้ดของโมเดลแบ่งส่วนถอดรหัสเป็น neck
ซึ่งทำหน้าที่ประกอบและรวมคุณลักษณะ กับ head ซึ่งให้ค่าความลึกออกมา ในรายงานนี้คำว่า DPT head
หมายถึงทั้งสองส่วน เช่นเดียวกับในบทนำ

\mbox{DA-v2} เองก็ฝึกด้วยครู (teacher) กล่าวคือ ฝึกโมเดลขนาดใหญ่ด้วยภาพสังเคราะห์ซึ่งมีความลึกที่แม่นยำก่อน
จากนั้นให้โมเดลนี้กำกับความลึกให้ภาพจริงที่ไม่มีความลึกราว 62 ล้านภาพ แล้วจึงให้โมเดลที่เผยแพร่เรียนรู้จากค่ากำกับที่โมเดลสร้างขึ้นนี้
ซึ่งเรียกว่า pseudo-label~\cite{da2} ผู้พัฒนา DINOv2 รายงานว่าทำนายความลึกได้ดีแม้ตรึงส่วนเข้ารหัสไว้
และฝึกเฉพาะส่วนถอดรหัส DPT ที่ต่อด้านบน~\cite{dinov2} การนำโมเดลที่ฝึกกับงานหรือโดเมนหนึ่งมาปรับใช้กับอีกงานหรือโดเมนหนึ่ง
เช่นนี้เรียกว่าการเรียนรู้แบบถ่ายโอน (transfer learning)~\cite{yosinski}
การฝึกเฉพาะส่วนถอดรหัสแบบที่งานนี้ทำใช้หน่วยความจำน้อยกว่ามาก
และยังคงคุณลักษณะทั่วไปที่ส่วนเข้ารหัสเรียนรู้มาจากภาพหลายล้านภาพไว้

\subsection{การฝึกด้วยความลึกจากเซนเซอร์}
โมเดลความลึกแบบมีผู้สอน (supervised) ใช้ความลึกที่วัดจากเซนเซอร์เป็นข้อมูลสอนมานานแล้ว โดยไม่ต้องให้คนกำกับ
ตัวอย่างเช่น NYU Depth~v2 ให้ความลึกในอาคารจากกล้อง Microsoft Kinect~\cite{nyuv2} KITTI ให้ความลึกบนถนนจากเครื่องสแกนเลเซอร์ที่ติดบนรถยนต์~\cite{kitti}
และ Eigen และคณะ~\cite{eigen14} ฝึกด้วยทั้งสองชุด สิ่งที่ต่างกันระหว่างงานเหล่านี้คือคุณภาพของเซนเซอร์
บน iPhone~Pro และ iPad~Pro นั้น ARKit รวมค่าที่วัดได้จาก LiDAR กับภาพสีเป็นแผนที่ความลึกขนาด
$256\times192$ พิกเซล พร้อมค่าความเชื่อมั่นของทุกพิกเซล~\cite{arkitscenes}
งานที่ทดสอบ LiDAR นี้เทียบกับวิธีสำรวจที่เป็นที่ยอมรับพบว่า แม่นยำราว 1~ซม. กับวัตถุที่ใหญ่กว่า 10~ซม.~\cite{luetzenburg}
และเพียงพอสำหรับการทำแผนที่สถาปัตยกรรมอย่างรวดเร็วที่มาตราส่วน 1:200~\cite{spreafico}
จึงมีเหตุผลที่จะใช้เป็นครู ARKitScenes~\cite{arkitscenes} ให้ความลึกนี้คู่กับความลึกจากเครื่องสแกนเลเซอร์ระดับงานสำรวจ
Faro~\cite{faro} ของห้องเดียวกัน โดยจัดให้ตรงกับเฟรมกล้องเดียวกัน และใช้ความลึกจาก Faro เป็นค่าอ้างอิง
ในโจทย์ที่เพิ่มความละเอียดของความลึกจาก LiDAR โดยอาศัยภาพสี
Prompt Depth Anything~\cite{promptda} ซึ่งเป็นงานที่ใช้ข้อมูลใกล้เคียงกับงานนี้ที่สุด ก็ป้อนความลึกจาก LiDAR ของ iPhone
เข้าโมเดล Depth Anything เป็นข้อมูลนำเข้าเพิ่มเติมเช่นกัน ในทั้งสองงาน LiDAR ต้องมีอยู่ทุกครั้งที่โมเดลทำงาน
ส่วนงานนี้ใช้ LiDAR เฉพาะตอนฝึก โมเดลที่นำไปใช้จึงต้องการเพียงภาพสี คำถามว่าครูที่มีสัญญาณรบกวนมากกว่านี้
ช่วยโมเดลขนาดใหญ่ที่ผ่านการฝึกล่วงหน้าได้เท่ากับเครื่องสแกนเลเซอร์หรือไม่
เมื่อตัดสินทั้งคู่ด้วยค่าอ้างอิงจากเลเซอร์ที่เป็นอิสระ คือเรื่องของวัตถุประสงค์ข้อ (1) และ (2)

\subsection{จากแผนที่ความลึกสู่แบบจำลองห้อง}
การรวมเฟรมต้องรู้ท่าทางกล้อง (camera pose) ของแต่ละเฟรม คือตำแหน่งและทิศทางของกล้อง
ARKit ติดตามท่าทางกล้องด้วย visual-inertial odometry (VIO) ซึ่งติดตามจุดเด่นในภาพจากเฟรมหนึ่งไปอีกเฟรมหนึ่ง
แล้วรวมกับเซนเซอร์วัดการเคลื่อนที่ของอุปกรณ์~\cite{vinsmono} เนื่องจากเซนเซอร์วัดการเคลื่อนที่วัดความเร่งเป็นหน่วยเมตริก
จุดที่ติดตามได้จึงมีตำแหน่งเป็นเมตร จุดเบาบาง (sparse points) เหล่านี้จึงให้สเกลและค่าเลื่อนที่โมเดลความลึกสัมพัทธ์ขาดไปได้ทีละเฟรม
แนวคิดการใช้ความลึกจากโครงข่ายร่วมกับเรขาคณิตแบบเบาบางของระบบติดตามมีมาตั้งแต่ CNN-SLAM~\cite{cnnslam}
และ MonoSDF~\cite{monosdf} ก็หาสเกลและค่าเลื่อนของทุกภาพก่อนนำความลึกจากภาพเดี่ยวไปช่วยสร้างพื้นผิว
การปรับสเกลด้วยจุดเบาบางของงานนี้ในหัวข้อ~\ref{sec:align} ก็อยู่ในแนวทางนี้

ขั้นต่อมาคือการสร้างห้อง ซึ่งทำโดยรวมความลึกเป็น truncated signed distance function (TSDF)~\cite{curless96}
โดยแบ่งพื้นที่เป็นลูกบาศก์เล็ก ๆ เรียกว่า voxel แต่ละ voxel เก็บระยะถึงพื้นผิวที่ใกล้ที่สุดที่สังเกตได้
เป็นบวกเมื่ออยู่หน้าพื้นผิวและเป็นลบเมื่ออยู่หลังพื้นผิว ตัดทิ้ง (truncate) เมื่อไกลเกินไม่กี่ voxel
แล้วเฉลี่ยค่านี้จากทุกเฟรมที่มองเห็น voxel นั้น จากนั้นจึงสกัดพื้นผิวตรงตำแหน่งที่ค่าเฉลี่ยเปลี่ยนเครื่องหมาย
ด้วยอัลกอริทึม marching cubes~\cite{marchingcubes} KinectFusion~\cite{kinectfusion} ทำให้การรวมแบบนี้ทำงานได้แบบเรียลไทม์
กับความลึกจาก Kinect และงานนี้ใช้การทำงานที่มีอยู่ใน Open3D~\cite{open3d}
การเฉลี่ยช่วยลดสัญญาณรบกวนที่สุ่มเปลี่ยนไปในแต่ละเฟรม แต่เฟรมที่มีสเกลไม่ตรงกันจะวางผนังเดียวกันไว้ที่ระยะต่างกัน
ค่าเฉลี่ยจึงทำให้ผนังเบลอหรือหนาขึ้น ความคงที่ของสเกลระหว่างเฟรม ซึ่งตัวชี้วัดรายภาพไม่ได้วัด จึงกำหนดคุณภาพของห้อง
นอกจากนี้ การรวมแบบ TSDF ไม่ขึ้นกับว่าความลึกเป็นเมตรนั้นมาจากแหล่งใด
จึงเป็นขั้นตอนคงที่ที่ยุติธรรมกับแหล่งความลึกทุกแหล่ง

ส่วนระบบสร้างแบบจำลองที่เรียนรู้ทั้งกระบวนการ (learned reconstruction) จะแทนที่บางขั้นตอนข้างต้น
ด้วยโครงข่ายประสาทเทียม Atlas~\cite{atlas} ทำนาย TSDF โดยตรง จากคุณลักษณะของภาพหลายภาพที่รู้ท่าทางกล้อง
NeuralRecon~\cite{neuralrecon}
ทำแบบเดียวกันทีละส่วนและแบบเรียลไทม์ และ SimpleRecon~\cite{simplerecon} ประมาณความลึกจากหลายเฟรมข้างเคียง
ก่อนรวมด้วยวิธีดั้งเดิม ทุกงานประเมินผลบน ScanNet~\cite{scannet} ซึ่งเป็นชุดการสแกนภาพสีและความลึกในอาคารขนาดใหญ่
แต่ละงานต้องใช้หลายเฟรมต่อการทำนายหนึ่งครั้ง และใช้โครงข่ายที่ฝึกมาสำหรับกระบวนการของตนเอง
จึงตอบคำถามต่างจากงานนี้ ซึ่งถามว่าโมเดลจากภาพเดี่ยวให้ห้องแบบใด เมื่อส่วนอื่นของกระบวนการไม่เปลี่ยน
งานนี้จึงไม่เปรียบเทียบกับงานกลุ่มนี้

\subsection{การวัดคุณภาพของแบบจำลองห้อง}
แบบจำลองห้องที่สร้างได้จะเทียบกับพื้นผิวอ้างอิงผ่านระยะระหว่างจุดที่สุ่มบนพื้นผิวทั้งสอง
Chamfer distance~\cite{chamfer} เฉลี่ยระยะจากแต่ละจุดถึงจุดที่ใกล้ที่สุดของอีกพื้นผิวทั้งสองทิศทาง
ส่วน Tanks and Temples~\cite{tanks} กำหนดระยะเกณฑ์ไว้ แล้วรายงานความแม่นยำ (precision)
คือสัดส่วนของพื้นผิวที่สร้างได้ที่อยู่ภายในระยะเกณฑ์จากพื้นผิวอ้างอิง ความครบถ้วน (recall)
คือสัดส่วนของพื้นผิวอ้างอิงที่อยู่ภายในระยะเกณฑ์จากพื้นผิวที่สร้างได้ และค่าเฉลี่ยฮาร์มอนิกของทั้งสอง คือ F-score
ระยะเกณฑ์ทำให้คะแนนผูกกับการใช้งาน ซึ่งเป็นวิธีที่บทนำใช้เชื่อมระยะ 2, 5 และ 10~ซม.
กับงานวัดเพื่อปรับปรุงห้อง งานวางเฟอร์นิเจอร์ และงานดูภาพรวมของห้อง สำหรับการวัดรายเฟรม
งานนี้ใช้ AbsRel และ \wone{} จากงานด้านการประมาณความลึกจากภาพเดี่ยว~\cite{eigen14}

\subsection{ตำแหน่งของงานวิจัยนี้}
โดยสรุป งานก่อนหน้ารายงานความแม่นยำรายภาพของโมเดลความลึกบนชุดทดสอบมาตรฐาน
ปรับโมเดลเมตริกด้วยข้อมูลกำกับจากภาพสังเคราะห์หรือจากเซนเซอร์ของชุดข้อมูลสาธารณะ
ป้อน LiDAR ของอุปกรณ์พกพาเข้าโมเดลขณะใช้งาน หรือทดสอบตัว LiDAR ของอุปกรณ์พกพาเองเทียบกับวิธีสำรวจ
เราไม่พบงานที่รวมประเด็นเหล่านี้ไว้แบบงานนี้ กล่าวคือ ใช้ LiDAR ของอุปกรณ์ผู้บริโภคเป็นครูเฉพาะตอนฝึก
เปรียบเทียบกับการใช้เครื่องสแกนเลเซอร์ระดับงานสำรวจเป็นครูบนห้องชุดเดียวกัน
และตัดสินทั้งคู่ด้วยค่าอ้างอิงจากเลเซอร์ที่เป็นอิสระ ทั้งเป็นเซนติเมตรบนแบบจำลองห้องและแบบรายเฟรม
